%% ============================================================
%% Appendix A - Source code
%% The listings are non-floating, following the pattern used by the
%% GUT template: a float cannot break across pages, so a long source
%% file placed inside one would be silently truncated.
%% ============================================================
\chapter{Source code}\label{app:source_code}

This appendix reproduces the parts of the implementation that are essential for understanding the system described in Chapter~\ref{chap:implementation}: the data structures exchanged between layers, one complete signal agent, the decision agent, and the portfolio simulator. The remaining modules, in particular the experiment harness and the machine-learning baselines, are not reproduced here; they are available in the project repository together with the data-preparation scripts and the result files.

The listings are verbatim copies of the source files, and the line numbering is that of the original files.

\section{Data structures}\label{app:data_structures}

The two dictionaries below are the entire interface between the layers of the system. A signal agent returns an \texttt{AgentSignal}; the decision agent consumes a list of them and returns a list of \texttt{Decision} records. Because both are declared as typed dictionaries, the contract is checkable statically while remaining a plain dictionary at runtime.

\begin{center}
  \captionsetup{type=listing}
  \captionof{listing}{The signal record returned by every signal agent.}
  \label{lst:agent_signal}
  \inputminted[linenos, frame=single, breaklines, fontsize=\footnotesize]{python}{other/AgentSignal.py}
\end{center}

\begin{center}
  \captionsetup{type=listing}
  \captionof{listing}{The decision record returned by the decision agent, including the per-agent attribution used in Section~\ref{subsec:practical}.}
  \label{lst:decision}
  \inputminted[linenos, frame=single, breaklines, fontsize=\footnotesize]{python}{other/Decision.py}
\end{center}

\section{A signal agent}\label{app:signal_agent}

Listing~\ref{lst:macd_agent} shows the MACD agent in full. It is representative of all five: a class attribute declaring the methodological family, a constructor taking the indicator parameters, and a single method that computes the indicator, maps it to a direction and a confidence, and returns the packed record. The neutral return in the guard clause is the abstention mechanism described in Section~\ref{subsec:agent_contract}.

\begin{center}
  \captionsetup{type=listing}
  \captionof{listing}{The MACD signal agent, an implementation of the momentum-family rule of Section~\ref{sec:agents}.}
  \label{lst:macd_agent}
  \inputminted[linenos, frame=single, breaklines, fontsize=\footnotesize]{python}{other/MACDAgent.py}
\end{center}

\section{The decision agent}\label{app:decision_agent}

Listing~\ref{lst:decision_agent} reproduces the decision agent in full. The family weights of \eqref{eq:family_weights} appear as the module-level default, and the mode-dependent thresholds of \eqref{eq:decision_rule} as the dictionary built in the constructor.

\begin{center}
  \captionsetup{type=listing}
  \captionof{listing}{The decision agent, implementing the confidence-weighted score and the thresholding rule.}
  \label{lst:decision_agent}
  \inputminted[linenos, frame=single, breaklines, fontsize=\footnotesize]{python}{other/DecisionAgent.py}
\end{center}

\section{The portfolio simulator}\label{app:portfolio_simulator}

Listing~\ref{lst:portfolio_simulator} reproduces the portfolio simulator, which implements the transaction-cost model of Section~\ref{subsec:costs} and produces the equity curve on which every metric in Chapter~\ref{chap:evaluation} is computed.

\begin{center}
  \captionsetup{type=listing}
  \captionof{listing}{The portfolio simulator: cash and position accounting net of commission, slippage and the one-way buy tax.}
  \label{lst:portfolio_simulator}
  \inputminted[linenos, frame=single, breaklines, fontsize=\footnotesize]{python}{other/PortfolioSimulator.py}
\end{center}
